%% pc-echelle-decibels — échelle des niveaux sonores L (dB) et intensités acoustiques I (W.m^-2).
%% L'axe des décibels est régulier ; en regard, les intensités progressent par puissances de 10
%% (c'est tout l'intérêt de l'échelle logarithmique).
%% Le dégradé vert -> orange -> rouge est obtenu par une suite de tranches de 5 dB (SVG léger).
\documentclass[tikz,border=4pt]{standalone}
\usepackage{liaisons}
\begin{document}
\begin{tikzpicture}[x=1cm,y=1cm,
  txt/.style={font=\scriptsize},
  ent/.style={font=\scriptsize\bfseries}]

%% échelle : 1 dB -> 0,0415 cm  (0 à 130 dB -> 5,4 cm) ; règle de x = 0 à x = 0,75
\def\ech{0.0415}
\def\rg{0.75}

%% ---------------- règle graduée en dégradé --------------------------------
\foreach \k in {0,1,...,12} {                     % 0 -> 65 dB : vert vers orange
  \pgfmathsetmacro\p{100-100*(5*\k+2.5)/65}
  \fill[solideE!\p!solideD] (0,{5*\k*\ech}) rectangle (\rg,{(5*\k+5)*\ech}); }
\foreach \k in {13,14,...,25} {                   % 65 -> 130 dB : orange vers rouge
  \pgfmathsetmacro\p{100-100*(5*\k+2.5-65)/65}
  \fill[solideD!\p!solideA] (0,{5*\k*\ech}) rectangle (\rg,{(5*\k+5)*\ech}); }
\draw[line width=0.6pt] (0,0) rectangle (\rg,{130*\ech});
\foreach \L in {0,10,...,130} {                   % graduations tous les 10 dB
  \draw[line width=0.5pt] (0,{\L*\ech}) -- (-0.09,{\L*\ech}); }

%% ---------------- en-têtes des colonnes -----------------------------------
\node[ent, anchor=east] at (-0.30,{142*\ech}) {$I$ (W$\cdot$m$^{-2}$)};
\node[ent, anchor=south] at (0.375,{131.5*\ech}) {$L$ (dB)};

%% ---------------- repères : niveau, intensité, situation -------------------
\foreach \L/\I/\lab in {%
    0/{$10^{-12}$}/{seuil d'audibilité ($I_0$)},
    40/{$10^{-8}$}/{chuchotement},
    60/{$10^{-6}$}/{conversation normale},
    80/{$10^{-4}$}/{conversation à haute voix},
    100/{$10^{-2}$}/{lecteur MP3},
    110/{$10^{-1}$}/{discothèque, concert},
    130/{$10^{1}$}/{réacteur d'avion\\\textbf{seuil de la douleur}}} {
  \draw[line width=0.6pt] (0,{\L*\ech}) -- (-0.16,{\L*\ech});
  \draw[line width=0.6pt] (\rg,{\L*\ech}) -- (\rg+0.09,{\L*\ech});
  \node[txt, anchor=east] at (-1.05,{\L*\ech}) {\I};
  \node[txt, anchor=east] at (-0.22,{\L*\ech}) {\L};
  \node[txt, anchor=west, align=left] at (\rg+0.14,{\L*\ech}) {\lab}; }

%% ---------------- encadré : définition du niveau sonore --------------------
\draw[line width=0.6pt, rounded corners=2pt] (1.95,0.18) rectangle (5.42,1.45);
\node[txt, anchor=center] at (3.68,1.27) {$L = 10\log\!\left(I/I_0\right)$};
\node[txt, anchor=center] at (3.68,1.00) {avec $I_0 = 1{,}0 \times 10^{-12}$ W$\cdot$m$^{-2}$};
\foreach \yy/\lab in {0.71/{$\times 2$ sur $I$ $\rightarrow$ $+3$ dB},
                      0.42/{$\times 10$ sur $I$ $\rightarrow$ $+10$ dB}} {
  \draw[solideD, line width=0.8pt, -{Stealth[length=4pt]}]
       (2.20,\yy-0.07) arc (200:-20:0.11);
  \node[txt, anchor=west, text=solideD] at (2.52,\yy) {\lab}; }
\end{tikzpicture}
\end{document}
